\documentclass[12pt]{article}

\usepackage{sbc-template}
\usepackage[utf8]{inputenc}
\usepackage[T1]{fontenc}
\usepackage[brazil]{babel}
\usepackage{graphicx}
\usepackage{amsmath,amssymb}
\usepackage{booktabs}
\usepackage{array}
\usepackage{url}
\usepackage{xcolor}
\usepackage{listings}
\usepackage{tikz}
\usetikzlibrary{positioning,arrows.meta,shapes.geometric}
\usepackage[hidelinks]{hyperref}
\usepackage{microtype}
\setlength{\emergencystretch}{3em}

\lstset{
  basicstyle=\ttfamily\small,
  breaklines=true,
  columns=fullflexible,
  keepspaces=true,
  frame=single,
  framesep=4pt,
  xleftmargin=4pt,
  numbers=none,
  showstringspaces=false
}

\newcolumntype{L}[1]{>{\raggedright\arraybackslash}p{#1}}

\title{Music Search Engine: um sistema de recuperação de informação para
música brasileira combinando modelos léxicos (TF-IDF e BM25) e busca
vetorial densa}

\author{Carsio Eddyo \and Carlos Alexandre \and Raquel de Sá \and Lelson
Nascimento}

\institute{Instituto de Computação (IComp) -- Universidade Federal do Amazonas
(UFAM)\\
Manaus -- AM -- Brasil\\
Disciplina ICC222 -- Tópicos em Recuperação de Informação -- 2026/1\\
\email{\{carsio, carlos.alexandre, raquel.sa, lelson\}@icomp.ufam.edu.br}}

\begin{document}

\maketitle

\begin{resumo}
Este trabalho descreve o \textit{Music Search Engine}, um sistema de
recuperação de informação (RI) para um catálogo curado de 50.000 faixas
musicais brasileiras, das quais 36.017 possuem letra consolidada. O sistema
implementa, a partir do zero, o núcleo clássico de RI --- pré-processamento
textual (normalização, tokenização, remoção de \textit{stopwords} e
\textit{stemming} RSLP/Snowball), índice invertido multi-campo e os modelos de
ranqueamento TF-IDF (com similaridade de cosseno) e Okapi BM25 --- e os
complementa com busca vetorial densa baseada em \textit{embeddings}. O motor é
exposto por uma API FastAPI consumida por uma SPA React, com classificação de
intenção determinística que roteia cada consulta para o índice adequado.
Apresentamos a metodologia de construção do corpus, as decisões de projeto e
uma avaliação experimental de desempenho: o índice é construído em cerca de
84~s sobre 36~mil documentos e responde consultas com latência mediana de
22,8~ms (BM25) e 17,9~ms (TF-IDF). Discutimos os \textit{trade-offs} entre os
modelos e delineamos a avaliação de relevância (Precision, MAP, nDCG) como
trabalho futuro.
\end{resumo}

\begin{abstract}
This work presents \textit{Music Search Engine}, an information retrieval (IR)
system over a curated catalog of 50,000 Brazilian music tracks, 36,017 of which
include consolidated lyrics. The system implements the classic IR core from
scratch --- text preprocessing (normalization, tokenization, stopword removal
and RSLP/Snowball stemming), a multi-field inverted index, and the TF-IDF
(cosine) and Okapi BM25 ranking models --- complemented by dense vector search
based on embeddings. The engine is exposed through a FastAPI backend consumed
by a React SPA, with a deterministic intent classifier that routes each query
to the appropriate index. We describe the corpus construction methodology,
design decisions, and a performance evaluation: the index is built in about
84~s over 36k documents and answers queries with a median latency of 22.8~ms
(BM25) and 17.9~ms (TF-IDF). We discuss model trade-offs and outline
relevance evaluation (Precision, MAP, nDCG) as future work.
\end{abstract}

% =====================================================================
\section{Introdução}
% =====================================================================

A Recuperação de Informação (RI) estuda como localizar, em uma coleção de
documentos, aqueles que melhor satisfazem uma necessidade de informação
expressa por uma consulta em linguagem natural~\cite{manning2008introduction}.
Sistemas de busca textual são uma das aplicações mais difundidas da
computação, e seus algoritmos fundamentais --- índice invertido, modelo
vetorial e funções de ranqueamento probabilístico --- permanecem no centro de
motores de busca modernos.

Este trabalho aplica esses conceitos, estudados na disciplina ICC222 (UFAM,
2026/1), à construção de um protótipo funcional de busca para um domínio
específico: a \textbf{música popular brasileira}. A escolha do domínio é
motivada por três características que o tornam um bom estudo de caso para RI:
(i)~as consultas são heterogêneas --- o usuário pode buscar por trecho de
\emph{letra}, por \emph{artista}, por \emph{gênero} ou por \emph{álbum}; (ii)~o
texto das letras é ruidoso (gírias, repetições, refrãos), o que estressa a
normalização e o ranqueamento; e (iii)~o vocabulário é majoritariamente em
português, exigindo \textit{stemming} adequado à língua.

Seguindo a Opção~1 da especificação da disciplina --- desenvolver um
\emph{protótipo funcional de sistema de recuperação de informação} ---,
construímos o \textit{Music Search Engine}: um sistema que implementa o núcleo
clássico de RI do zero e o disponibiliza por meio de uma aplicação web. As
principais contribuições deste trabalho são:

\begin{itemize}
  \item a curadoria de um corpus de 50.000 faixas brasileiras, enriquecido com
    36.017 letras consolidadas a partir de múltiplas fontes públicas;
  \item a implementação, em Python, de um núcleo de RI auditável contendo
    pré-processamento, índice invertido multi-campo e os modelos TF-IDF e
    Okapi BM25, com \textit{field boosts} por domínio;
  \item a integração de um modelo vetorial denso (\textit{embeddings} +
    banco vetorial Milvus) como caminho complementar para consultas
    semânticas;
  \item uma aplicação web completa (API FastAPI + SPA React) com classificação
    de intenção determinística e roteamento por tipo de entidade; e
  \item uma avaliação experimental de desempenho (tempo de indexação, latência
    de consulta, tamanho do índice) com números reais e reproduzíveis.
\end{itemize}

O restante do texto está organizado da seguinte forma. A
Seção~\ref{sec:fundamentacao} revisa os fundamentos de RI utilizados. A
Seção~\ref{sec:metodologia} detalha a metodologia: construção do corpus,
núcleo de RI, motores de busca e arquitetura do sistema. A
Seção~\ref{sec:resultados} apresenta os resultados experimentais. A
Seção~\ref{sec:discussao} discute os achados e as limitações, e a
Seção~\ref{sec:conclusao} conclui e aponta trabalhos futuros.

% =====================================================================
\section{Fundamentação teórica}
\label{sec:fundamentacao}
% =====================================================================

\subsection{Pré-processamento e índice invertido}

Antes de qualquer comparação, o texto é convertido em uma sequência de
\textbf{termos} normalizados. O pipeline canônico aplica, em ordem:
normalização (\textit{lowercasing}, remoção de acentos via NFKD, substituição
de pontuação por espaços), tokenização, remoção de \textit{stopwords} e
\textit{stemming}. O \textit{stemming} reduz flexões a um radical comum
(p.~ex.\ ``cantando'', ``cantamos'' $\rightarrow$ ``cant''), de modo que
variações morfológicas casem com a consulta. Para o português utiliza-se o
algoritmo RSLP~\cite{orengo2001stemming} e, para o inglês, o
Snowball~\cite{bird2009nltk}.

O \textbf{índice invertido}~\cite{manning2008introduction} é a estrutura que
torna a busca eficiente: para cada termo, armazena-se a \textit{postings list}
--- a lista de documentos em que ele ocorre e com qual frequência. Formalmente,
mantém-se $\text{postings}[t] = [(d, \text{tf}(t,d)), \dots]$, além do
comprimento de cada documento e do mapeamento entre identificadores internos e
externos. Assim, responder ``quais documentos contêm o termo $t$?'' custa
$O(1)$, e o ranqueamento percorre apenas os documentos candidatos.

\subsection{Modelo vetorial e TF-IDF}

No modelo vetorial~\cite{salton1975vector}, documentos e consultas são vetores
em um espaço cujas dimensões são os termos do vocabulário. O peso de um termo
$t$ em um documento $d$ combina sua frequência local e sua raridade global:
%
\begin{equation}
  w(t,d) = \text{tf\_weight}(\text{tf}(t,d)) \cdot \text{idf}(t),
  \qquad
  \text{idf}(t) = \ln\!\left(\frac{N}{\text{df}(t)}\right),
\end{equation}
%
em que $N$ é o número total de documentos e $\text{df}(t)$ o número de
documentos que contêm $t$~\cite{sparckjones1972statistical}. Termos raros
recebem maior peso. A função $\text{tf\_weight}$ admite três variantes:
\emph{raw} ($\text{tf}$ direto), \emph{log} ($1 + \ln \text{tf}$, o padrão
adotado) e \emph{augmented} ($0{,}5 + 0{,}5\cdot\text{tf}/\max_t \text{tf}$). A
relevância é medida pela similaridade de cosseno entre os vetores:
%
\begin{equation}
  \text{cos\_sim}(d,q)
    = \frac{\mathbf{v}_d \cdot \mathbf{v}_q}
           {\lVert \mathbf{v}_d \rVert \, \lVert \mathbf{v}_q \rVert}.
\end{equation}

\subsection{Okapi BM25}

O BM25~\cite{robertson2009probabilistic} substitui o cosseno por um escore
aditivo que incorpora duas correções importantes: a \textbf{saturação} da
frequência de termos (repetir um termo 50 vezes não vale 50 vezes mais) e a
\textbf{normalização pelo tamanho} do documento. O escore é:
%
\begin{equation}
  \text{BM25}(D,Q)
    = \sum_{t \in Q} \text{IDF}_{\text{bm25}}(t) \cdot
      \frac{\text{tf}(t,D)\,(k_1 + 1)}
           {\text{tf}(t,D) + k_1\!\left(1 - b + b\,\dfrac{|D|}{\text{avgdl}}\right)},
\end{equation}
%
com o IDF suavizado:
%
\begin{equation}
  \text{IDF}_{\text{bm25}}(t)
    = \ln\!\left(\frac{N - \text{df}(t) + 0{,}5}{\text{df}(t) + 0{,}5} + 1\right).
\end{equation}
%
O parâmetro $k_1$ controla a velocidade de saturação da $\text{tf}$ e $b$ pondera a
normalização por tamanho. Adotam-se os valores convencionais de
Lucene/Elasticsearch~\cite{lucene_similarity}, $k_1 = 1{,}5$ e $b = 0{,}75$.

\subsection{Busca vetorial densa}

Modelos léxicos casam termos exatos e, por isso, falham diante de sinônimos e
paráfrases. A \textbf{busca vetorial densa} representa documentos e consultas
como vetores contínuos (\textit{embeddings}) produzidos por redes neurais
\cite{reimers2019sentence}, recuperando por similaridade semântica. A consulta
``música triste'' pode então recuperar faixas com termos como ``melancolia''
ou ``lamento'', mesmo sem sobreposição léxica. A busca por vizinhos mais
próximos sobre os vetores é viabilizada por bancos vetoriais como o
Milvus~\cite{wang2021milvus}.

\subsection{Métricas de avaliação de RI}

A qualidade de um sistema de RI é tipicamente medida com
\textit{relevance judgments}: dado um conjunto de consultas e os documentos
relevantes para cada uma, computam-se \textbf{Precision@k} (fração de
relevantes entre os $k$ primeiros), \textbf{Recall@k}, \textbf{Mean Average
Precision} (MAP) e \textbf{nDCG} (que pondera ganho por
posição)~\cite{manning2008introduction}. Essas métricas requerem um
\emph{golden set} anotado --- elemento ainda não construído neste trabalho e
discutido na Seção~\ref{sec:conclusao}.

% =====================================================================
\section{Metodologia}
\label{sec:metodologia}
% =====================================================================

O sistema é organizado em camadas, isolando os algoritmos puros de RI do
domínio musical e da apresentação. A Figura~\ref{fig:arquitetura} resume a
arquitetura.

\begin{figure}[ht]
\centering
\begin{tikzpicture}[
    font=\small,
    box/.style={draw, rounded corners, align=center, minimum height=0.9cm,
                inner sep=4pt, fill=blue!5},
    core/.style={box, fill=orange!12},
    arr/.style={-{Stealth[length=2mm]}, thick}
  ]
  \node[box, text width=10.5cm] (front) {Frontend SPA (React + Vite) \quad / \quad UIs Tkinter (debug)};
  \node[box, below=0.5cm of front, text width=10.5cm] (api)
    {API FastAPI \;--\; classificação de intenção (heurística) + roteamento};
  \node[box, below=0.5cm of api, text width=5.0cm] (motors)
    {Motores \\ \texttt{SparseSearchEngine} \\ \texttt{MultiEntityIndex}};
  \node[box, right=0.5cm of motors, text width=4.8cm] (vector)
    {Busca vetorial \\ \textit{embeddings} + Milvus};
  \node[core, below=0.5cm of motors, text width=5.0cm] (core)
    {Núcleo de RI \\ preprocessing $\rightarrow$ índice invertido $\rightarrow$ ranking (TF-IDF, BM25)};
  \node[box, below=0.5cm of core, text width=10.5cm] (data)
    {Corpus curado \;\texttt{br\_curated\_lyrics.parquet}\; + dimensões de entidades};

  \draw[arr] (front) -- (api);
  \draw[arr] (api) -- (motors);
  \draw[arr] (api.south) to[out=-90,in=90] (vector.north);
  \draw[arr] (motors) -- (core);
  \draw[arr] (core) -- (data);
\end{tikzpicture}
\caption{Arquitetura em camadas do \textit{Music Search Engine}. O núcleo de RI
(laranja) não conhece o domínio musical; os motores aplicam o núcleo ao corpus
curado e adicionam \textit{field boosts} e roteamento.}
\label{fig:arquitetura}
\end{figure}

\subsection{Construção do corpus}

A fonte primária é o \textit{Spotify Metadata}, com metadados estruturados de
milhões de faixas (nome, artistas, álbum, gêneros, popularidade,
\textit{audio features}). Uma etapa de análise exploratória e curadoria filtra
o conjunto para faixas brasileiras (gêneros contendo ``brazilian'', ``samba'',
``bossa nova'' etc.; artistas com mercado BR; popularidade mínima), produzindo
50.000 faixas.

Os metadados não bastam: para indexar letras é preciso o \emph{texto}. Ele é
obtido por uma \textbf{cascata de fontes públicas} (LRCLib, Lyrics.ovh,
Vagalume, Letras.mus.br, Genius e LyricFind): tenta-se a primeira fonte e, em
caso de ausência, passa-se à seguinte, com variantes de título (remoção de
sufixos como ``\textit{Remastered}'') e \textit{retry} com \textit{backoff}
exponencial. Toda a coleta é memoizada em SQLite e protegida por
\textit{rate limiting} (\textit{token bucket}) e \textit{circuit breaker},
tornando o processo idempotente e resiliente. Ao final, 36.017 faixas obtêm
letra consolidada.

Para responder consultas por artista e gênero com painéis informativos, um
pipeline de \textbf{enriquecimento determinístico} coleta texto da
Wikipédia~PT e o materializa em dimensões separadas. O dataset final é
versionado em tabelas Parquet, acompanhado de um manifesto
(\texttt{br\_dataset\_manifest.json}, versão 0.3.0) com contagens, tamanhos e
\textit{hashes} SHA-1. A Tabela~\ref{tab:dataset} resume o \textit{snapshot}
atual.

\begin{table}[ht]
\centering
\caption{Snapshot versionado do dataset (manifesto v0.3.0).}
\label{tab:dataset}
\begin{tabular}{lrl}
\toprule
\textbf{Artefato} & \textbf{Registros} & \textbf{Uso} \\
\midrule
\texttt{br\_curated\_tracks} & 50.000 & metadados estruturados das faixas \\
\texttt{br\_curated\_lyrics} & 36.017 & corpus textual indexado (BM25/TF-IDF) \\
\texttt{br\_artists}         & 7.255  & dimensão de artistas (Wikipédia PT) \\
\texttt{br\_genres}          & 42     & dimensão de gêneros \\
\bottomrule
\end{tabular}
\end{table}

\subsection{Núcleo de RI}

O núcleo (camada \texttt{core/}) é deliberadamente agnóstico ao domínio:
recebe documentos e consultas genéricos e devolve escores. O
pré-processamento (\texttt{core/preprocessing.py}) implementa o pipeline da
Seção~\ref{sec:fundamentacao}. O índice (\texttt{core/indexer.py}) é
\textbf{multi-campo}: cada documento possui campos (\texttt{track\_name},
\texttt{lyrics}, \texttt{artist\_names}, \texttt{artist\_genres},
\texttt{macro\_genre}, \texttt{album\_name}) e o índice mantém \textit{postings}
separados por campo, com identificadores internos densos e persistência via
\textit{pickle}. O ranqueamento (\texttt{core/ranking.py}) implementa TF-IDF e
BM25 sobre o mesmo índice, pré-calculando IDFs e normas via \textit{caching}; o
custo de uma consulta é $O(|Q| \cdot |\text{candidatos}|)$.

\subsection{Motores de busca e \textit{field boosts}}

{\sloppy
A camada de motores (\texttt{motors/}) aplica o núcleo ao domínio. O
\texttt{SparseSearchEngine} calcula o escore de cada campo separadamente,
normaliza-o por campo (para combinar escalas distintas) e soma ponderado por
pesos de domínio (\textit{field boosts}), apresentados na
Tabela~\ref{tab:weights}. A lógica é: um termo da consulta que aparece na
\emph{letra} ou no \emph{título} é sinal mais forte de relevância do que no
nome do álbum.\par}

\begin{table}[ht]
\centering
\caption{Pesos por campo (\texttt{DEFAULT\_FIELD\_WEIGHTS}).}
\label{tab:weights}
\begin{tabular}{lcl}
\toprule
\textbf{Campo} & \textbf{Peso} & \textbf{Justificativa} \\
\midrule
\texttt{lyrics}        & 4{,}0  & conteúdo principal \\
\texttt{track\_name}   & 2{,}5  & título é forte sinal de relevância \\
\texttt{artist\_genres}& 1{,}5  & gênero ajuda a filtrar \\
\texttt{artist\_names} & 1{,}0  & nome do artista \\
\texttt{macro\_genre}  & 1{,}0  & macro-categoria curada \\
\texttt{album\_name}   & 0{,}75 & sinal mais fraco \\
\bottomrule
\end{tabular}
\end{table}

O \texttt{MultiEntityIndex} é uma fachada que conhece várias entidades
(tracks, artists, albums, genres, composers) e expõe
\texttt{search\_routed(query, intent, top\_k)}, despachando a consulta ao índice
adequado. Quando uma dimensão ainda não existe no \textit{snapshot} (p.~ex.\
álbuns), o roteamento degrada graciosamente para o índice de tracks.

\subsection{Busca vetorial densa}

Como caminho complementar e opcional, o sistema indexa as faixas como
\textit{embeddings} (via Ollama com \texttt{nomic-embed-text}, 768 dimensões,
ou OpenAI \texttt{text-embedding-3-small}, 1536 dimensões), persistidos em
Milvus Lite, e responde por similaridade vetorial. Esse motor não está no
caminho crítico da API --- é exposto por CLI e interface Tkinter
dedicadas --- mas demonstra a recuperação semântica complementar aos modelos
léxicos.

\subsection{Classificação de intenção e arquitetura web}

A consulta ``saudade do meu amor'' tem cara de letra; ``ana carolina'', de
artista. Um classificador \textbf{heurístico determinístico} em
\texttt{web/app.py} inspeciona marcadores (prefixos como ``album'', frases
longas com pronomes, nomes próprios curtos, gêneros conhecidos) e atribui uma
intenção, que orienta o roteamento. A escolha por heurística, e não por um
LLM, é deliberada: é instantânea, não depende de chave nem de
\textit{cold start}, e o critério é auditável em código --- uma chamada externa
síncrona no caminho de cada consulta seria um gargalo grave no servidor de
produção.

O motor é exposto por uma API \textbf{FastAPI}~\cite{ramirez2018fastapi} que,
na inicialização, carrega o \texttt{SparseSearchEngine} e o
\texttt{MultiEntityIndex}. Os principais \textit{endpoints} são
\texttt{/api/search} (busca roteada), \texttt{/api/search/lyric} (busca em
letras com \textit{snippets} numerados e destaque dos termos),
\texttt{/api/artist/\{id\}}, \texttt{/api/album/\{id\}} e
\texttt{/api/song/\{id\}}. O \textit{frontend} é uma SPA em React~+~Vite~+~
TypeScript que consome essa API.

\subsection{Protocolo experimental}
\label{sec:protocolo}

A avaliação de desempenho foi conduzida por um \textit{script} reproduzível
(\texttt{docs/report/benchmark.py}) que consome o \texttt{SparseSearchEngine}
--- sem reimplementar lógica de ranqueamento. Mediram-se: (i)~o tempo de
construção a frio (\textit{cold build}) do índice invertido sobre as 36.017
letras e o tamanho do artefato persistido; (ii)~o vocabulário por campo; e
(iii)~a latência de consulta para BM25 e TF-IDF sobre um conjunto fixo de 18
consultas representativas dos três regimes (letra, artista, gênero), com 3
consultas de aquecimento descartadas e 5 repetições por consulta (90 amostras
por algoritmo). Os experimentos rodaram em uma máquina de desenvolvimento
(Apple Silicon, macOS), em ambiente Python 3.12 gerenciado por \texttt{uv}.

% =====================================================================
\section{Resultados}
\label{sec:resultados}
% =====================================================================

\subsection{Indexação e tamanho do índice}

A construção a frio do índice invertido multi-campo sobre as 36.017 faixas com
letra levou \textbf{84,3~s}, e o artefato persistido (\textit{pickle}) ocupa
\textbf{24,9~MB} em disco. Esse índice é gerado uma única vez e reaproveitado
nas execuções seguintes, de modo que o custo não recai sobre a consulta. O
tempo é dominado pelo pré-processamento --- em especial o \textit{stemming}
RSLP de seis campos textuais, sendo o campo \texttt{lyrics} o mais volumoso. A
Tabela~\ref{tab:vocab} mostra o vocabulário (termos distintos) por campo: as
letras concentram a maior parte do vocabulário (61.821 termos), coerente com
seu papel de conteúdo principal.

\begin{table}[ht]
\centering
\caption{Vocabulário (termos distintos após pré-processamento) por campo.}
\label{tab:vocab}
\begin{tabular}{lr}
\toprule
\textbf{Campo} & \textbf{Termos distintos} \\
\midrule
\texttt{lyrics}        & 61.821 \\
\texttt{track\_name}   & 9.809 \\
\texttt{artist\_names} & 9.212 \\
\texttt{album\_name}   & 8.995 \\
\texttt{artist\_genres}& 355 \\
\texttt{macro\_genre}  & 19 \\
\bottomrule
\end{tabular}
\end{table}

\subsection{Latência de consulta}

A Tabela~\ref{tab:latency} apresenta a latência de ponta a ponta da função de
busca (incluindo pré-processamento da consulta, \textit{lookup} no índice,
ranqueamento multi-campo e composição dos resultados). Ambos os algoritmos
respondem em poucas dezenas de milissegundos: a mediana é de 22,8~ms para BM25
e 17,9~ms para TF-IDF, com p95 abaixo de 90~ms. A diferença entre os dois é
pequena e esperada, dado que compartilham o mesmo índice e estrutura de
varredura; o BM25 é marginalmente mais custoso por avaliar a normalização por
tamanho de documento a cada termo.

\begin{table}[ht]
\centering
\caption{Latência de consulta (ms), 90 amostras por algoritmo.}
\label{tab:latency}
\begin{tabular}{lrrrrr}
\toprule
\textbf{Algoritmo} & \textbf{Média} & \textbf{Mediana} & \textbf{p95} & \textbf{Mín} & \textbf{Máx} \\
\midrule
BM25   & 34{,}4 & 22{,}8 & 89{,}2 & 0{,}5 & 106{,}3 \\
TF-IDF & 30{,}0 & 17{,}9 & 79{,}7 & 0{,}4 & 97{,}4 \\
\bottomrule
\end{tabular}
\end{table}

\subsection{Comparação qualitativa BM25 $\times$ TF-IDF}

A Tabela~\ref{tab:qualitative} compara o \textit{top}-3 dos dois modelos em
consultas representativas. O caso ``chega de saudade'' é ilustrativo: o BM25
ranqueia corretamente o clássico ``Chega de Saudade'' (João Gilberto) nas
primeiras posições, enquanto o TF-IDF, por não saturar a frequência do termo
``saudade'', favorece faixas que apenas repetem esse termo no título. Para
``bossa nova'', cujos resultados dependem mais de correspondência exata, ambos
concordam. Isso é coerente com a literatura: o BM25 tende a oferecer
ranqueamento de melhor qualidade em corpora com documentos de tamanhos
variados, como letras de música~\cite{robertson2009probabilistic}.

\begin{table}[ht]
\centering
\caption{Top-3 por modelo (consultas selecionadas). Artista entre parênteses.}
\label{tab:qualitative}
\renewcommand{\arraystretch}{1.15}
\begin{tabular}{c L{4.9cm} L{4.9cm}}
\toprule
\textbf{\#} & \textbf{BM25} & \textbf{TF-IDF} \\
\midrule
\multicolumn{3}{l}{\textit{Consulta:} \texttt{chega de saudade}} \\
\midrule
1 & Chega de Saudade (João Gilberto) & Saudade de Você (Claudia Leitte) \\
2 & Chega de Saudade (João Gilberto) & Saudade Que Fala -- Ao Vivo (Pagode do Adame) \\
3 & Saudade (Fi Barreto) & Saudade de Você (Filhos De Jorge) \\
\midrule
\multicolumn{3}{l}{\textit{Consulta:} \texttt{bossa nova}} \\
\midrule
1 & bossa nova (aupinard) & bossa nova (aupinard) \\
2 & Nova Bossa Nova (Marcos Valle) & Nova Bossa Nova (Marcos Valle) \\
3 & Rio da Bossa Nova (Ambulante Discos) & Rio da Bossa Nova (Ambulante Discos) \\
\bottomrule
\end{tabular}
\end{table}

% =====================================================================
\section{Discussão}
\label{sec:discussao}
% =====================================================================

Os resultados confirmam que o protótipo atende ao objetivo de busca textual
eficiente: latências de dezenas de milissegundos viabilizam uma experiência
interativa, e a construção do índice, embora custosa (84~s), ocorre uma única
vez fora do caminho da consulta.

A comparação qualitativa evidencia o principal \textit{trade-off} entre os
modelos léxicos. O \textbf{TF-IDF} é simples e rápido, mas vulnerável à
dominância de termos repetidos. O \textbf{BM25}, pela saturação da
$\text{tf}$ e normalização por tamanho, produz ranqueamento mais robusto ---
razão pela qual é o padrão da indústria e o algoritmo \textit{default} do
sistema. A \textbf{busca vetorial densa} cobre um regime ortogonal (consultas
semânticas e paráfrases), porém ao custo de gerar e armazenar
\textit{embeddings} e de maior latência; por isso permanece como caminho
complementar, fora do caminho crítico da API.

A \textbf{classificação de intenção heurística} é eficaz e barata, mas tem
limites: consultas ambíguas (um nome que também é palavra comum) podem ser
roteadas incorretamente. A arquitetura mitiga isso com \textit{fallback} para
o índice de tracks, mas uma classificação mais sofisticada poderia melhorar a
precisão do roteamento sem reintroduzir dependência de serviços externos.

As principais \textbf{limitações} do trabalho são: (i)~a cobertura de letras
(36.017 de 50.000 faixas, $\approx$72\%), limitada pela disponibilidade nas
fontes públicas; (ii)~as dimensões de álbuns e compositores ainda não
enriquecidas no \textit{snapshot} atual; e, sobretudo, (iii)~a ausência de uma
\textbf{avaliação de relevância} quantitativa. Sem um \textit{golden set}
anotado não é possível reportar Precision@k, MAP ou nDCG --- a avaliação
apresentada limita-se a desempenho e a inspeção qualitativa. Essa é a lacuna
mais relevante a ser fechada e o foco natural do trabalho futuro.

% =====================================================================
\section{Conclusão e trabalhos futuros}
\label{sec:conclusao}
% =====================================================================

Este trabalho apresentou o \textit{Music Search Engine}, um protótipo
funcional de RI para música brasileira que implementa, do zero, o núcleo
clássico da disciplina --- pré-processamento, índice invertido multi-campo,
TF-IDF e BM25 --- e o disponibiliza por meio de uma aplicação web com busca
vetorial densa complementar. Construímos um corpus curado de 50.000 faixas
(36.017 com letra) e demonstramos, com números reais, que o sistema indexa
36~mil documentos em $\approx$84~s e responde consultas com latência mediana
inferior a 23~ms, com o BM25 entregando ranqueamento qualitativamente superior
ao TF-IDF em consultas de letra.

Como \textbf{trabalho futuro}, destacam-se: (i)~construir um \textit{golden
set} de 20--50 consultas com julgamentos de relevância no formato
\textit{qrels} e implementar o módulo de avaliação (Precision@k, Recall@k,
MAP, nDCG), hoje um \textit{placeholder} no código; (ii)~conduzir um
\textit{benchmark} comparativo controlado entre BM25, TF-IDF e busca vetorial
em termos de qualidade; (iii)~concluir o enriquecimento das dimensões de
álbuns e compositores; e (iv)~explorar abordagens \textit{híbridas} que
combinem o escore léxico e o semântico. Esses passos transformariam a
avaliação qualitativa atual em uma análise quantitativa rigorosa, completando o
ciclo metodológico da disciplina.

% =====================================================================
\bibliographystyle{plain}
\bibliography{refs}
% =====================================================================

\end{document}
